\subsection{The next skeleton cannot use only differential-generated one-forms}

\begin{proposition}[A-S09-23: Proved; a precise obstruction to the free zero-form second skeleton]
\label{AS0923}
Let $N\geq2$ and let $P_{\leq2}$ be a two-truncated,
nonnegative weight-graded strict PD cochain $S$-algebra
whose first skeleton is exactly $P(N)$ of (SK.1).
Assume its weight-zero zero-form algebra is the constant
coefficient ring.  Fix a weight $w$ for which the actual
special-differential map is onto:
\[
 \Omega^1_{\mathrm{PD}}(P_2^{*,0}/S)^w
                  \longrightarrow P_2^{w,1},
                  \qquad a\,db\longmapsto a\,d_Pb.
 \tag{S09.108}
\]
For every odd prime $p>3(N-1)$ for which (S09.108) is
onto in weight $w=p$, the row-one homology of this
truncated object satisfies
\[
                 H_1\bigl(N(P^{p,1})\bigr)
                         \otimes\mathbb Z_{(p)}\ne0.
 \tag{S09.109}
\]
Here $H_1$ uses only simplex degrees zero, one and two.
Consequently (S09.108) cannot be onto in every sufficiently
large weight for the first two skeletons of an admissible
theory.  In particular, choosing $P_2$ to be a free
even-zero-form PD disk algebra does not give the eventual
row-one exactness required for a second-skeleton version
of A-004.

No cellularity assumption is made on the given $P_2$.
The argument constructs the cellular object and the
actual map from this cellular cover needed for the cotangent comparison.
The restriction $N\geq2$ is essential to the cited
cotangent calculation; no assertion for $N=1$ is made.
\end{proposition}

\begin{proof}
Write $D=P^{*,0}$ for the zero-form algebra.  Begin with
the cellular simplicial PD $S$-algebra generated by the
specified one-skeleton: its degree-one generators are
the $x_w$, $w\geq N$, with endpoints $0,T^{[w]}$.
Its higher terms initially have only the forced
degenerate generators.  For every positive-weight
element
\[
                         z\in N_2(D)
                           =\ker d_1\cap\ker d_2
\]
adjoin a free PD two-cell whose normalized top generator
$u_z$ has boundary $d_0z$.  This is an allowed cell
because the simplicial identities give
$d_0(d_0z)=d_1(d_0z)=0$.
More explicitly, the corresponding free simplicial
PD pushout is on the inclusion
$K(1)\hookrightarrow D(2)$, with the bottom generator
sent to $d_0z$; its quotient parameter is $K(2)$.
Every cell therefore leaves simplex degrees zero and
one unchanged.  Denote the resulting cellular algebra
by $Q^{(2)}$.

There is an actual truncated PD algebra map
\[
                  f:Q^{(2)}_{\leq2}\longrightarrow D_{\leq2}
                  \quad\hbox{with}\quad f(u_z)=z,
 \tag{S09.110}
\]
which is the identity in degrees zero and one.
Its degree-two component is surjective.  To verify this
directly, the normalization projector in degree two is
\[
 \pi_2=(1-s_0d_1)(1-s_1d_2),\qquad
 a=\pi_2a+s_1d_2a+s_0(d_1a-d_2a).
 \tag{S09.111}
\]
The two faces $d_1,d_2$ kill $\pi_2a$, by the simplicial
identities.  The first term is the image of its indexed
generator $u_{\pi_2a}$, with zero used when $\pi_2a=0$;
the other two terms are images of known degeneracies
of $D_1=Q^{(2)}_1$.  This is an explicit preimage of
every homogeneous $a\in D_2$.  Weight zero is already
the coefficient algebra by assumption.

The universal property of special differentials gives
an actual map of truncated simplicial modules
\[
      \Omega^1_{\mathrm{PD}}(Q^{(2)}/S)_{\leq2}
                        \longrightarrow P^{*,1}_{\leq2}.
 \tag{S09.112}
\]
It is an isomorphism in simplex degrees zero and one:
both degree-zero terms are zero, and $P(N)_1$ is the
free PD de Rham algebra on its zero-form generators.
It is surjective in degree two in every weight where
(S09.108) is onto.  Indeed write a one-form as a sum
of $a\,d_Pb$ and lift both zero-forms by (S09.111).
Applying the natural normalization projector proves
surjectivity on $N_2$ as well.
Since the maps on $N_0,N_1$ are isomorphisms, the
images of $N_2\to N_1$ on the two sides are equal.
Thus (S09.112) induces an isomorphism on $H_1$ in
every weight in question.  This step supplies the
comparison with a potentially noncellular $P_2$;
it is not a replacement made without a map.

The zero-form component of $Q^{(2)}$ has
\[
                             \pi_0Q^{(2)}=S/I_N.
 \tag{S09.113}
\]
In fact the degree-one endpoint differences contain
every $T^{[w]}$ with $w\geq N$.  Conversely the two
endpoint maps agree after quotienting by $I_N$, so
their difference ideal is precisely $I_N$.
Higher-dimensional cells leave these degrees unchanged.

Complete $Q^{(2)}$ to an actual cellular resolution
$Q\to S/I_N$ using only free cells of simplicial
dimensions at least two.  One direct construction is
to adjoin, at each stage, a cell bounding every
positive-dimensional normalized cycle in the current
algebra and then repeat.  Each cycle specifies the
map from its sphere parameter, and the free disk
pushout supplies its stated boundary.  The augmentation
sends each new top generator to zero.  In the filtered
union every positive-dimensional cycle occurs in a
finite stage and is a boundary at the next stage.
The zero and first simplex algebras never change, so
(S09.113) remains fixed.  This is the required actual
cellular resolution, with its inclusion from
$Q^{(2)}$.

That inclusion induces an actual special-differential
map which is the identity in simplex degrees zero
and one.  Hence
\[
 H_1\Omega^1_{\mathrm{PD}}(Q^{(2)}/S)^w
               \longrightarrow H_1\Omega^1_{\mathrm{PD}}(Q/S)^w
                  \quad\hbox{is onto}.             \tag{S09.114}
\]
Equivalently new two-cells can add boundaries in this
homological degree, while higher cells do not change
it.  No invariance under the underived free PD functor
is used.

Apply A-S09-12 with $k=1$ to the displayed cellular
resolution.  For every odd prime $p>3(N-1)$ its
right-hand group in (S09.114), in weight $p$, is
nonzero after localization at $p$.  The surjection
(S09.114) and the proved isomorphism (S09.112) give
(S09.109).  Such primes are unbounded, which contradicts
eventual exactness of the form-one row whenever
(S09.108) is eventually onto.  A free even-zero-form
disk algebra has (S09.108) an isomorphism, so it is
covered by this conclusion.
\end{proof}
